\chapter{Thực nghiệm và Đánh giá}

\section{Cấu hình thực nghiệm}
\begin{table}[H]
\centering
\begin{tabularx}{\textwidth}{|l|X|}
\hline
\textbf{Thuộc tính} & \textbf{Giá trị} \\
\hline
Tên tập dữ liệu & CD Covers Dataset (TMBuD) \\
\hline
Tổng số ảnh trong CSDL & 99 ảnh \\
\hline
Số ảnh truy vấn (Test) & 10 ảnh \\
\hline
Giá trị K & 3, 5, 11, 21 \\
\hline
Phần cứng & CPU Tiêu chuẩn \\
\hline
Cách xác định ảnh liên quan & Ảnh có chung tiền tố ID (ví dụ \texttt{001\_x.png}). Ground truth được định nghĩa qua tên file. \\
\hline
\end{tabularx}
\caption{Cấu hình hệ thống thực nghiệm.}
\end{table}

\section{Bộ tham số của các phương pháp}
\begin{table}[H]
\centering
\begin{tabularx}{\textwidth}{|l|X|l|}
\hline
\textbf{Phương pháp} & \textbf{Tham số chính} & \textbf{Độ đo} \\
\hline
Color Histogram & Không gian HSV, 32 bins/kênh & Chi-Squared \\
\hline
Color Correlogram & 8 màu lượng tử hóa, 4 khoảng cách & Chi-Squared \\
\hline
SIFT & 500 từ vựng (BoVW), Max 500 keypoints & Cosine \\
\hline
ORB & 500 từ vựng (BoVW), Max 500 keypoints & Euclidean \\
\hline
HOG & $64 \times 64$ resize, block 16, cell 8, 9 bins & Cosine \\
\hline
LBP & Bán kính 1, 8 điểm lân cận, 256 bins & Chi-Squared \\
\hline
Color + HOG & Nối Vector, chuẩn hóa L2 Norm & Cosine \\
\hline
\end{tabularx}
\caption{Bảng siêu tham số của từng phương pháp.}
\end{table}

\section{Kết quả MAP (Mean Average Precision)}
Bảng dưới đây trình bày điểm số đánh giá độ chính xác trung bình (MAP) do module \texttt{evaluate\_map} tự động trích xuất dựa trên 10 ảnh Test.

\begin{table}[H]
\centering
\begin{tabularx}{\textwidth}{|l|c|c|c|c|c|}
\hline
\textbf{Phương pháp} & \textbf{Độ đo khoảng cách} & \textbf{MAP@3} & \textbf{MAP@5} & \textbf{MAP@11} & \textbf{MAP@21} \\
\hline
Color Histogram & Chi-Squared & 0.033 & 0.057 & \textbf{0.066} & 0.074 \\
\hline
Color Correlogram & Chi-Squared & 0.028 & 0.023 & 0.036 & 0.062 \\
\hline
SIFT (Inverted Index) & TF-IDF Cosine & 0.033 & 0.033 & 0.058 & 0.058 \\
\hline
ORB (Inverted Index) & TF-IDF Cosine & 0.033 & 0.033 & 0.033 & 0.035 \\
\hline
HOG & Cosine & 0.011 & 0.011 & 0.020 & 0.031 \\
\hline
LBP & Chi-Squared & 0.033 & 0.048 & 0.059 & 0.067 \\
\hline
Color\_HOG & Cosine & 0.028 & 0.043 & 0.046 & 0.054 \\
\hline
\end{tabularx}
\caption{Bảng kết quả MAP trên 10 ảnh Test sau khi tích hợp chỉ mục ngược.}
\end{table}

\section{Phân tích kết quả MAP}
Dựa trên số liệu thực nghiệm sau khi tích hợp chỉ mục Inverted Index, có thể rút ra một số nhận xét quan trọng:
\begin{itemize}
    \item \textbf{Sự thay đổi MAP của SIFT và ORB:} So với phương pháp so khớp brute-force L2/Cosine nguyên bản trên dense histogram, việc tích hợp Inverted Index kết hợp trọng số TF-IDF làm thay đổi nhẹ chỉ số MAP (SIFT giảm xuống 0.058 ở MAP@21 và ORB xuống 0.035). Nguyên nhân là do mô hình TF-IDF tập trung đánh giá mức độ tương quan thông qua độ thưa (sparsity) và độ hiếm của visual words thay vì khoảng cách hình học L2 thuần túy. Tuy nhiên, sự đánh đổi này đem lại hiệu quả vượt trội về mặt tốc độ tìm kiếm.
    \item \textbf{Color Histogram thể hiện bất ngờ:} Trái với nhận định đặc trưng màu là thô sơ, Color Histogram (MAP@21 = 0.074) đạt kết quả rất tốt trên tập đĩa CD. Điều này đến từ việc các bìa đĩa CD thường mang một dải màu đặc trưng và dễ phân biệt.
    \item \textbf{HOG hoạt động cực kỳ kém:} Với số điểm MAP@21 chỉ đạt 0.031, chứng tỏ dải phân bố Gradient không đủ để mô tả độ chi tiết trong ảnh bìa CD vốn mang nhiều phông chữ và texture lộn xộn.
    \item \textbf{LBP ở mức trung bình khá:} Với MAP@21 = 0.067, LBP bắt được kết cấu chữ in tốt hơn nhiều so với HOG.
    \item \textbf{Sự kết hợp Color\_HOG chưa tạo ra đột biến:} Điểm MAP ở ngưỡng 0.054 cho thấy việc nối vector HOG vào Color đã vô tình làm giảm độ chính xác của Color gốc. Cần áp dụng trọng số $\alpha$ và $\beta$ hợp lý hơn thay vì ghép tỉ lệ 1:1.
\end{itemize}

\section{Đánh giá thời gian truy vấn (Query Latency)}
Để chứng minh hiệu quả vượt trội về mặt hiệu năng của cấu trúc Inverted Index, hệ thống thực hiện đo đạc thời gian so khớp thực tế (Matching Latency) trên tập dữ liệu CD Covers (TMBuD) với quy mô mở rộng là $1.363$ ảnh.

Thực nghiệm đo thời gian so khớp trung bình trên 200 lượt chạy cho một ảnh query bằng đặc trưng SIFT:
\begin{itemize}
    \item \textbf{Brute-Force (Cosine Matching):} $0.834$ ms.
    \item \textbf{Inverted Index (TF-IDF Matching):} $0.385$ ms.
    \item \textbf{Tỷ lệ tăng tốc (Speedup Factor):} $2.17\times$ (nhanh hơn gấp đôi).
\end{itemize}
\textbf{Nhận xét:} Dù CSDL thực nghiệm tương đối nhỏ ($1.363$ ảnh) và đặc trưng BoW SIFT có mật độ phân bố khá dày (hơn 650 ngàn posting entries), cấu trúc Inverted Index đã cho thấy ưu thế rõ rệt khi chỉ thực hiện đối chiếu trên các visual word có mặt trong query histogram. Sự tăng tốc này sẽ càng rõ rệt (dạng lũy thừa tuyến tính) khi số lượng ảnh trong cơ sở dữ liệu nâng cấp lên hàng trăm nghìn hoặc hàng triệu ảnh.


\section{Minh họa kết quả thực nghiệm}
Dưới đây là một số ví dụ minh họa kết quả khi thử nghiệm truy vấn với từng loại đặc trưng trên giao diện ứng dụng thực tế.

\subsection{Color Histogram}
\missingfigure{color_histogram_01.png}{Kết quả truy vấn Color Histogram – trường hợp 1.}
\missingfigure{color_histogram_02.png}{Kết quả truy vấn Color Histogram – trường hợp 2.}
\missingfigure{color_histogram_03.png}{Kết quả truy vấn Color Histogram – trường hợp 3.}

\subsection{Color Correlogram}
\missingfigure{color_correlogram_01.png}{Kết quả Color Correlogram 1.}
\missingfigure{color_correlogram_02.png}{Kết quả Color Correlogram 2.}
\missingfigure{color_correlogram_03.png}{Kết quả Color Correlogram 3.}

\subsection{SIFT}
\missingfigure{sift_01.png}{Kết quả SIFT 1.}
\missingfigure{sift_02.png}{Kết quả SIFT 2.}
\missingfigure{sift_03.png}{Kết quả SIFT 3.}

\subsection{ORB}
\missingfigure{orb_01.png}{Kết quả ORB 1.}
\missingfigure{orb_02.png}{Kết quả ORB 2.}
\missingfigure{orb_03.png}{Kết quả ORB 3.}

\subsection{HOG}
\missingfigure{hog_01.png}{Kết quả HOG 1.}
\missingfigure{hog_02.png}{Kết quả HOG 2.}
\missingfigure{hog_03.png}{Kết quả HOG 3.}

\subsection{LBP}
\missingfigure{lbp_01.png}{Kết quả LBP 1.}
\missingfigure{lbp_02.png}{Kết quả LBP 2.}
\missingfigure{lbp_03.png}{Kết quả LBP 3.}

\subsection{Color + HOG}
\missingfigure{color_hog_01.png}{Kết quả Color + HOG 1.}
\missingfigure{color_hog_02.png}{Kết quả Color + HOG 2.}
\missingfigure{color_hog_03.png}{Kết quả Color + HOG 3.}
